% Generated draft: requires mathematical and source review.
\documentclass[11pt]{article}
\usepackage[margin=25mm]{geometry}
\usepackage{fontspec}
\setmainfont{Latin Modern Roman}
\usepackage{amsmath,amssymb,mathtools,unicode-math}
\setmathfont{Latin Modern Math}
\usepackage{enumitem,needspace,xurl,hyperref}
\setlength{\parindent}{0pt}
\setlength{\parskip}{5pt}
\setlength{\emergencystretch}{3em}
\newcommand{\N}{\mathbb{N}}
\newcommand{\Z}{\mathbb{Z}}
\newcommand{\Q}{\mathbb{Q}}
\newcommand{\R}{\mathbb{R}}
\newcommand{\C}{\mathbb{C}}
\newcommand{\F}{\mathbb{F}}
\newcommand{\A}{\mathbb{A}}
\newcommand{\PP}{\mathbb{P}}
\newcommand{\E}{\mathbb{E}}
\newcommand{\eps}{\varepsilon}
\newcommand{\dd}{\,\mathrm{d}}
\begin{document}
\section*{UnsolvedMath ID: 20001042}
\textbf{Source problem code:} catalogue entry
\subsection{Definitions and mathematical statement}
The literal source fragment is only ``Conjecture 6.5. (Seymour) Under the hypotheses of'', with its antecedent and conclusion missing. The verification calls it severely truncated despite the registry's \texttt{exact} label. The following load conjecture is the attached aid's reconstruction, not a complete formula visible in the original extract.

For a finite loopless simple digraph $D$ on $n$ vertices, assign each arc a subset $L(e)$ of a finite label set; empty-list arcs may be ignored. Define labelled outdegree/load
\[
 \ell(v)=\sum_{e\text{ leaving }v}|L(e)|=\sum_i d^+_{E_i}(v),
 \qquad E_i=\{e:i\in L(e)\}.
\]
A rainbow directed cycle chooses a distinct available label for every one of its arcs. Digons count as cycles.

\textbf{Reconstructed Seymour load conjecture.} If $\ell(v)\geq n$ for every vertex, there is a rainbow directed cycle. The missing original hypotheses must be verified before treating this reconstruction as an exact source statement.

\textbf{Reported sharp two-label theorem.} If all lists lie in a set of at most two labels and there is no rainbow directed cycle, the aid proves
\[
 \sum_{v\in V(D)}\ell(v)\leq n(n-1).
\]
Thus average load greater than $n-1$ forces a rainbow digon, and the reconstructed pointwise-$n$ conjecture holds in this case. The bound is attained even with both labels present: on a transitive tournament give every arc both labels, obtaining total load $2\binom n2=n(n-1)$ and no directed cycle, for $n\geq2$. The theorem concerns total/average load; the extremal example need not have all individual loads equal.

\textbf{Reported threshold obstruction.} For any integer $d\geq1$, a directed cycle on $d+1$ vertices with the same $d$ available labels on every arc has load $d$ at each vertex but no rainbow cycle: its only directed cycle is longer than the label set. In particular, replacing the pointwise threshold $n$ by $n-1$ would fail. The report attributes the need for this example to the missing surrounding source discussion, which is not supplied.

\textbf{Reported three-label frontier.} If all arcs have all three labels, then $\ell(v)=3d^+(v)$ and a rainbow cycle is exactly a directed cycle of length at most three. The resulting restriction is therefore
\[
 \min_v d^+(v)\geq\lceil n/3\rceil
 \quad\Longrightarrow\quad D\text{ has a cycle of length at most }3,
\]
the directed-triangle Caccetta--Haggkvist case (with loops excluded and digons allowed as short witnesses). The two-label proof does not establish this three-label case or the unrestricted load conjecture. The report records partial progress, not a general resolution.
\subsection{Short English statement}
The original conjecture is truncated. Under the aid's list-labelled reconstruction, a pointwise load of at least the vertex count should force a rainbow cycle. The aid proves the sharp two-label total-load bound, gives a threshold counterexample one unit lower, and identifies the all-three-label case with the directed-triangle Caccetta--Haggkvist problem.
\subsection{Sources}
\begin{itemize}
\item Local UnsolvedMath registry, permanent ID $20001042$, \texttt{catalogue entry}; canonical \texttt{aim-combinatorics-notes.json}, \texttt{notes[166]}, source item 6.
\item AIM workshop ``The Caccetta--Haggkvist conjecture,'' PDF fragment of Conjecture 6.5 attributed to Seymour: \url{https://aimath.org/WWN/caccetta/caccetta.pdf}. The canonical verification expressly says the record is severely truncated.
\item Research attempt 1, \texttt{partial\_result}, run status \texttt{open}, internal rigor \texttt{proved}, low novelty confidence; current registry partially solved. Explicitly machine-generated, not peer reviewed, and to be independently verified before citation. No complete external reference or original literature entry supplied.
\end{itemize}
\end{document}
